% ============================================================
%  Documento IEEE — Conjunto Reductor
%  Compilar en Overleaf con compilador pdfLaTeX
% ============================================================
\documentclass[conference]{IEEEtran}
\IEEEoverridecommandlockouts

% ── Paquetes esenciales ──────────────────────────────────────
\usepackage[utf8]{inputenc}
\usepackage[T1]{fontenc}
\usepackage[spanish]{babel}
\usepackage{graphicx}
\usepackage{float}
\usepackage{booktabs}
\usepackage{array}
\usepackage{tabularx}
\usepackage{multirow}
\usepackage{amsmath}
\usepackage{amssymb}
\usepackage{hyperref}
\usepackage{url}
\usepackage{xurl}
\usepackage{cite}
\usepackage{caption}
\usepackage{colortbl}
\usepackage[table]{xcolor}

% ── Configuración español ────────────────────────────────────
\addto\captionsspanish{%
  \renewcommand{\tablename}{Tabla}%
  \renewcommand{\figurename}{Fig.}%
}

% ─────────────────────────────────────────────────────────────
\begin{document}

% ══════════════════════════════════════════════════════════════
%  TÍTULO Y AUTORES (IEEE)
% ══════════════════════════════════════════════════════════════
\title{Sistema de Engranajes para Banda Transportadora: Diseño,
Selección de Materiales y Análisis Mecánico}

\author{
  \IEEEauthorblockN{Miguel A. Rojas, Sebastián A. Rodríguez,
  David A. Rodriguez, Jose A. Rincon, Daniel García Araque}
  \IEEEauthorblockA{
    \textit{Ingeniería Mecatrónica}\\
    \textit{Universidad Militar Nueva Granada}\\
    Bogotá, Colombia\\
    \{miguel.arojas, sebastian.rodriguez, david.arodrigu1,
     jose.rincon, daniel.garcia\}@unimilitar.edu.co
  }
}

\maketitle

% ══════════════════════════════════════════════════════════════
%  ABSTRACT E INDEX TERMS
% ══════════════════════════════════════════════════════════════
\begin{abstract}
Se presenta el diseño de un conjunto reductor de velocidad orientado a su aplicación en
una banda transportadora de carga media, enfocado en la adaptación de velocidad y par en
sistemas de transmisión de potencia. El estudio aborda la configuración del sistema, el
análisis de sus componentes principales y la selección de materiales en función de sus
propiedades mecánicas y condiciones de operación. Asimismo, se consideran criterios de
diseño mecánico para garantizar un comportamiento adecuado bajo carga continua,
asegurando resistencia al desgaste, confiabilidad y eficiencia en el funcionamiento. El
modelo CAD desarrollado en SolidWorks incluye la verificación de restricciones de
ensamblaje y el análisis de interferencias entre componentes. Los factores de seguridad
adoptados se justifican explícitamente a partir de criterios de fatiga establecidos. El
resultado es una solución técnicamente viable para aplicaciones industriales de manejo
de materiales.
\end{abstract}

\begin{IEEEkeywords}
conjunto reductor, transmisión de potencia, banda transportadora, engranajes, diseño
mecánico, factor de seguridad
\end{IEEEkeywords}

% ══════════════════════════════════════════════════════════════
%  I. DEFINICIÓN DEL SISTEMA
% ══════════════════════════════════════════════════════════════
\section{Definición del Sistema}

\subsection{Descripción General}

El conjunto reductor es un mecanismo capaz de reducir la velocidad a la salida de este,
aumentando el par torsor y manteniendo la potencia y una relación de transmisión fija
\cite{tarancon2018}. En la actualidad, las velocidades que proporcionan los motores
eléctricos industriales son demasiado altas, y el rango de velocidades que ofrecen no son
las adecuadas para satisfacer ciertas aplicaciones; por ende, la gran mayoría de máquinas
accionadas por un motor necesitarán de estos mecanismos, con el fin de adaptar la
velocidad del motor a la requerida por la máquina.

\subsubsection{Funcionamiento}

El principio fundamental para entender el funcionamiento del conjunto reductor es la
transmisión de potencia. Los engranajes se encargan de transmitir potencia de un motor a
una carga, reduciendo la velocidad de rotación y aumentando el par. El eje de entrada está
acoplado al motor e impulsa un engranaje más pequeño, el cual va acoplado a un engranaje
más grande para generar una ventaja mecánica. Este fenómeno se conoce como principio de
relación de transmisión, donde la relación entre los engranajes conductor y conducido
determina la reducción de velocidad y la multiplicación del par. Una relación de 10:1
significa que el eje de salida gira a una décima parte de la velocidad, entregando
aproximadamente diez veces el par \cite{aox}.

\subsubsection{Aplicaciones}

Los conjuntos reductores se emplean en diversas aplicaciones industriales debido a su
capacidad de adaptar las condiciones de velocidad y par entre un motor y una carga. En
robótica, son utilizados para mejorar la precisión y el control del movimiento, permitiendo
reducir la velocidad de los actuadores y aumentar el torque disponible en las articulaciones.
En sistemas de ventilación industrial, los reductores permiten ajustar la velocidad de
operación de los ventiladores según las necesidades del proceso. Por otro lado, en equipos
de elevación como grúas y polipastos, son fundamentales para incrementar el par y permitir
el levantamiento seguro de cargas pesadas \cite{lopez2023}.

Sin embargo, una de las aplicaciones más representativas de estos sistemas se encuentra
en las bandas transportadoras industriales, donde el conjunto reductor transforma la alta
velocidad de rotación de un motor en una velocidad adecuada para el desplazamiento
controlado de materiales, proporcionando además el par necesario para el movimiento de la
carga \cite{conveyor}. En este tipo de sistemas, la caja de engranajes constituye el núcleo
del mecanismo de transmisión, ya que se encarga de transferir la potencia hacia elementos
como poleas, rodillos o cadenas. La confiabilidad del reductor es un factor crítico en la
operación de estos sistemas, debido a que los transportadores suelen trabajar bajo
condiciones de carga continua y durante largos periodos de tiempo. En industrias como la
minería, puertos y plantas de procesamiento de materiales, una falla en la caja reductora
puede generar interrupciones significativas en múltiples operaciones \cite{guomao2026}.

% ══════════════════════════════════════════════════════════════
%  II. ESPECIFICACIONES DE OPERACIÓN
% ══════════════════════════════════════════════════════════════
\section{Especificaciones de Operación}

\subsection{Variables Principales}

Para determinar variables como la potencia y la velocidad del mecanismo, se parte de la
premisa de que el sistema reductor está concebido para operar en una banda transportadora
de uso industrial. Su propósito es adaptar la velocidad y el par entregados por un motor
eléctrico a las condiciones específicas del proceso. En la práctica, los motores eléctricos
industriales suelen operar en un rango de velocidades comprendido entre 1500 y 1800\,rpm.
Con el fin de seleccionar la opción más adecuada en términos de relación costo-beneficio,
se evaluaron diferentes motores eléctricos industriales, cuyos resultados se presentan en
la Tabla~\ref{tab:comparacion-motores}.

\begin{table}[H]
\centering
\caption{Comparación de Motores Eléctricos Industriales}
\label{tab:comparacion-motores}
\begin{tabular}{lccccc}
\hline
\textbf{Modelo / marca} & \textbf{kW} & \textbf{rpm} & \textbf{A} &
\textbf{\%$\eta$} & \textbf{USD} \\
\hline
WEG W22 (IE2)        & 22    & 1760 & 35{,}1 & 93 & 3\,301 \\
Siemens SIMOTICS GP  & 15    & 1500 & 28     & 94 & 3\,088 \\
WEG W22 (IE3)        & 7{,}5 & 1500 & 13     & 90 & 550    \\
\hline
\multicolumn{6}{l}{\small\textit{Nota.} Datos obtenidos de hojas técnicas de fabricantes.}
\end{tabular}
\end{table}

A partir del análisis de las opciones presentadas en la Tabla~\ref{tab:comparacion-motores},
se seleccionó un motor asíncrono trifásico de jaula de ardilla de 15\,kW y 4 polos, con
una velocidad nominal aproximada de 1500\,rpm. Esta elección se justifica debido a que
motores de menor potencia, como los de 7{,}5\,kW, pueden resultar insuficientes para
aplicaciones de transporte continuo bajo carga, mientras que motores de mayor potencia,
como los de 22\,kW, implican un sobredimensionamiento del sistema, incrementando
innecesariamente el costo y el consumo energético. Por lo tanto, el motor de 15\,kW
representa un punto de equilibrio adecuado entre capacidad de carga, eficiencia energética
y costo. Adicionalmente, este tipo de motores se caracteriza por su alta confiabilidad, bajo
mantenimiento y capacidad de operar en servicio continuo (S1), lo que los hace apropiados
para aplicaciones industriales como bandas transportadoras.

\subsection{Tipo de Carga}

El sistema opera bajo condiciones de carga continua, ya que las bandas transportadoras
están diseñadas para funcionar durante largos periodos de tiempo sin interrupciones. Este
tipo de carga se clasifica como carga constante o ligeramente fluctuante, debido a que el
material transportado puede generar variaciones moderadas en el esfuerzo requerido.
Durante el arranque y la detención del sistema pueden presentarse cargas dinámicas; sin
embargo, estas no son predominantes durante la operación normal. Por esta razón, se
considera necesario un factor de servicio mayor a 1, el cual permite contemplar condiciones
de sobrecarga ocasionales. El factor de servicio indica la capacidad del motor para soportar
cargas superiores a su valor nominal. Por ejemplo, un factor de servicio de 1{,}25 implica
que el motor puede operar hasta con un 125\,\% de su carga nominal sin comprometer su
funcionamiento \cite{snellgroup}. El motor seleccionado cumple con este criterio, lo que
lo hace adecuado para este tipo de aplicación.

\subsection{Condiciones Ambientales}

El sistema reductor y el motor operan en condiciones propias de entornos industriales, lo
que implica exposición a factores como polvo, humedad y variaciones de temperatura. Estos
elementos pueden afectar el desempeño del sistema si no se consideran adecuadamente
durante el diseño. Por esta razón, es fundamental seleccionar materiales resistentes al
desgaste y a la corrosión, así como implementar sistemas de sellado que protejan los
componentes internos del reductor. Esto permite evitar la contaminación de los lubricantes
y reducir el desgaste de los engranajes. Adicionalmente, el uso de motores con grado de
protección IP55 o IP56 garantiza la protección contra el ingreso de partículas sólidas y
salpicaduras de agua, asegurando un funcionamiento confiable bajo condiciones industriales
exigentes.

% ══════════════════════════════════════════════════════════════
%  III. MODELO CAD
% ══════════════════════════════════════════════════════════════
\section{Modelo CAD}

El modelo CAD desarrollado en el software SolidWorks corresponde a un conjunto reductor
de engranajes, diseñado con el objetivo de representar la estructura y disposición de los
elementos que componen el sistema de transmisión. El modelo incluye componentes como
ejes, engranajes, rodamientos y carcasa, permitiendo visualizar la interacción entre estos
elementos durante la operación del sistema.

\begin{figure}[H]
  \centering
  \includegraphics[width=0.85\columnwidth]{Imagen1}
  \caption{Conjunto reductor, vista isométrica.}
  \label{fig:isometrica}
\end{figure}

En la Fig.~\ref{fig:isometrica} se presenta la vista isométrica del modelo CAD del conjunto
reductor, en la cual se observa el ensamblaje completo del sistema. Se identifican
claramente la carcasa principal de forma prismática, la base de soporte con pernos de
fijación y los ejes de entrada y salida ubicados en la cara frontal del conjunto.

\begin{figure}[H]
  \centering
  \includegraphics[width=0.85\columnwidth]{imagen2}
  \caption{Vista interna del sistema de engranajes.}
  \label{fig:interna}
\end{figure}

En la Fig.~\ref{fig:interna} se presenta la vista interna del conjunto reductor, donde se
ha retirado la carcasa con el fin de visualizar la disposición de los elementos de
transmisión. En esta representación se identifican claramente los engranajes montados sobre
ejes paralelos, así como los elementos de soporte y acoplamiento entre ellos.

\begin{figure} [H]
    \centering
    \includegraphics[width=1\linewidth]{Explosionado.png}
    \caption{Explosionado del sistema}
    \label{fig:placeholder1}
\end{figure}

Como se puede ver en la Fig.~\ref{fig:placeholder1}, rea realizó el explosionado del CAD donde se puede visualizar mejor la composición del sistema para la caja de velocidades y como se encuentran distribuidas las piezas internas en los ejes.

\subsection{Proceso de Modelado CAD: Restricciones y Análisis de Interferencias}

El modelo fue construido de forma jerárquica en SolidWorks, comenzando por el modelado
individual de cada pieza y procediendo al ensamblaje mediante la aplicación sistemática de
restricciones geométricas (\textit{mates}). A continuación se detallan los aspectos más
relevantes del proceso:

\subsubsection{Restricciones de ensamblaje}

Las principales restricciones aplicadas durante el proceso de ensamblaje fueron:

\begin{itemize}
  \item \textbf{Coincidencia de ejes}: los ejes del piñón y la rueda dentada se definieron
  como paralelos, con la distancia interaxial fijada en $c = r_1 + r_2 = 95 + 155 =
  250$\,mm, garantizando el engrane correcto entre los dientes.
  \item \textbf{Restricciones de posición axial}: se emplearon relaciones de plano
  coincidente para fijar la posición longitudinal de los engranajes sobre los ejes,
  evitando desplazamientos axiales no deseados que podrían generar desalineación.
  \item \textbf{Fijación de la carcasa}: la carcasa se restringió completamente (6 grados
  de libertad bloqueados) para que actuara como referencia fija de todo el ensamblaje.
  \item \textbf{Relaciones de rodamientos y tapetas}: las tapetas ciegas y pasantes fueron
  ensambladas con restricciones de concentricidad y coplanaridad respecto a los orificios
  de la carcasa, asegurando el alineamiento de los apoyos.
  \item \textbf{Chavetas y circlips}: se aplicaron restricciones de simetría y
  concentricidad para los elementos de sujeción sobre los ejes, verificando que no
  existiesen conflictos con los hombros o cambios de sección.
\end{itemize}

\subsubsection{Análisis de interferencias}

Previo a la generación de planos y a la validación analítica del diseño, se ejecutó la
herramienta de detección de interferencias (\textit{Interference Detection}) de SolidWorks
sobre el ensamblaje completo. Los resultados obtenidos fueron los siguientes:

\begin{itemize}
  \item \textbf{Sin interferencias volumétricas} entre las piezas en condición de operación
  estática. Ningún componente sólido presenta solapamiento con otro en la posición de
  montaje final.
  \item \textbf{Holgura mínima de 0{,}5\,mm} en la zona de contacto entre los dientes de
  los engranajes, coherente con el módulo $m = 10$\,mm y el juego de cabeza estándar
  ($c^* = 0{,}25$), lo que permite la circulación adecuada del lubricante entre los
  flancos.
  \item \textbf{Ajuste correcto de rodamientos}: los rodamientos 635-20312 fueron
  modelados con ajuste $H7/k6$ en el alojamiento de la carcasa y $k6/h6$ en el eje, lo
  que garantiza la interferencia de prensa requerida para la aplicación de carga radial
  continua, según lo recomendado por el catálogo de rodamientos SKF \cite{skf2018}.
  \item \textbf{Circlips y anillos de seguridad} verificados sin conflicto geométrico con
  los hombros del eje ni con las tapetas pasantes, asegurando que los desplazamientos
  axiales queden correctamente limitados.
\end{itemize}

Estos resultados confirman la viabilidad de ensamblaje del diseño antes de proceder a la
fabricación de los componentes, reduciendo el riesgo de errores dimensionales en la
manufactura.

\subsection{Justificación del Diseño para la Aplicación Seleccionada}

El modelo CAD diseñado para una banda transportadora de carga media presenta una
adecuada correspondencia con los requerimientos de esta aplicación. Su configuración
estructural permite una integración eficiente entre las condiciones de entrada del motor y
las exigencias de salida del sistema. La geometría del conjunto, basada en ejes paralelos y
engranajes rectos, facilita una transmisión de potencia eficiente y confiable. Esta
configuración es ampliamente utilizada en entornos industriales debido a su simplicidad
constructiva, facilidad de mantenimiento y buen desempeño bajo condiciones de carga
continua. Asimismo, la distribución de esfuerzos entre los engranajes y los ejes contribuye
a reducir concentraciones de carga, favoreciendo la durabilidad del sistema y disminuyendo
el desgaste de los componentes mecánicos.

Por otra parte, el uso de engranajes de diferentes diámetros permite una reducción
progresiva de la velocidad, lo cual resulta adecuado para sistemas que requieren un
movimiento controlado y constante, como las bandas transportadoras. El diseño incorpora
además una carcasa que protege el mecanismo frente a condiciones ambientales adversas,
como la presencia de polvo y humedad. Adicionalmente, con el fin de mitigar efectos
asociados a vibraciones durante la operación, se incluye una base que permite la fijación
del conjunto a superficies rígidas, mejorando la estabilidad del sistema. Finalmente, las
dimensiones del modelo han sido ajustadas para mantener coherencia con el motor
previamente seleccionado, garantizando así una adecuada integración entre los elementos
del sistema.

A continuación, se describen de manera detallada los componentes que conforman el modelo.

\subsection{Identificación de Elementos}

En la Fig.~\ref{fig:componentes} se presenta la identificación de los principales
componentes del conjunto reductor. A partir de esta numeración, en la
Tabla~\ref{tab:elementos-conjunto-reductor} se describen los elementos, su cantidad y su
función dentro del sistema.

\begin{figure}[H]
  \centering
  \includegraphics[width=0.85\columnwidth]{imagen3}
  \caption{Identificación de componentes del conjunto reductor. La figura muestra la
  disposición general de los elementos internos y externos del sistema, facilitando la
  identificación de cada pieza dentro del ensamble.}
  \label{fig:componentes}
\end{figure}

\begin{table*}[t]
\centering
\caption{Identificación de Elementos del Conjunto Reductor}
\label{tab:elementos-conjunto-reductor}
\small
\begin{tabular}{clcp{9cm}}
\hline
\textbf{N.\textsuperscript{o}} & \textbf{Elemento} & \textbf{Cant.} & \textbf{Función} \\
\hline
1  & Carcasa               & 1  & Protege los componentes internos y mantiene la alineación del sistema. \\
2  & Rueda dentada         & 1  & Transmite movimiento entre los ejes y contribuye a la reducción de velocidad. \\
3  & Eje                   & 2  & Soporta los engranajes y transmite el movimiento mecánico. \\
4  & Piñón                 & 1  & Recibe el movimiento de entrada y actúa como elemento conductor. \\
5  & Tapeta ciega          & 2  & Cierra y protege los extremos del conjunto. \\
6  & Tapeta pasante        & 2  & Permite el paso del eje y su soporte adecuado. \\
7  & Base soporte          & 1  & Sirve para fijar el conjunto a la estructura externa. \\
8  & Cuña 1.6x1x7  & 2  & Asegura la unión entre eje y elemento giratorio. \\
9  & Cuña 2x1.2x5  & 2  & Transmite el par sin deslizamiento entre piezas. \\
10 & Rodamiento 635-20312  & 4  & Reduce la fricción y soporta el giro de los ejes. \\
11 & DIN 1804 M85×2        & 4  & Elemento de sujeción y ajuste. \\
12 & Anillo Caucho 13.9-13   & 2  & Evita desplazamientos axiales. \\
13 & Anillo Caucho 7.7-6  & 2  & Ayuda al sellado y retención de lubricante. \\
14 & Circlip DIN 471 60×3  & 4  & Asegura componentes sobre el eje. \\
15 & DIN EN 24015 M12×60   & 16 & Tornillería de fijación. \\
16 & DIN EN 24015 M16×55   & 8  & Tornillería de fijación. \\
17 & DIN EN 24015 M20×65   & 4  & Tornillería de fijación. \\
18 & Tapón llenado   & 1  & Permite el llenado de lubricante. \\
19 & Visor nivel aceite  & 1 & Permite verificar el nivel de lubricación. \\
20 & Tapón desagüe  & 1  & Facilita el vaciado del lubricante. \\
\hline
\end{tabular}
\end{table*}

La identificación de estos componentes permite comprender la función de cada elemento
dentro del conjunto reductor y la forma en que contribuyen a la transmisión de potencia,
la sujeción estructural, la lubricación y la protección del sistema. Esta organización facilita
el análisis del diseño y su relación con la aplicación industrial seleccionada.

A partir de esta caracterización, es posible establecer criterios adecuados para la selección
de materiales, considerando las condiciones de operación, los esfuerzos mecánicos y el
entorno al que estará sometido cada componente, lo cual será abordado en la siguiente
sección.

% ══════════════════════════════════════════════════════════════
%  IV. DISEÑO DE EJE Y CÁLCULOS DEL SISTEMA
% ══════════════════════════════════════════════════════════════
\section{Diseño de Eje y Cálculos del Sistema}

El sistema consiste en un reductor de engranajes cilíndricos rectos con las siguientes
especificaciones de entrada:

\begin{table}[H]
\centering
\caption{Parámetros de Entrada del Motor}
\label{tab:parametros}
\begin{tabular}{lll}
\hline
\textbf{Parámetro} & \textbf{Símbolo} & \textbf{Valor} \\
\hline
Potencia                          & $P$   & $15\;\mathrm{kW}$   \\
Velocidad del motor (eje entrada) & $n_1$ & $1500\;\mathrm{RPM}$ \\
Módulo de engranaje               & $m$   & $10\;\mathrm{mm}$   \\
Número de dientes --- Piñón       & $Z_1$ & $19$                \\
Número de dientes --- Rueda       & $Z_2$ & $31$                \\
\hline
\end{tabular}
\end{table}

\subsection{Relación de Transmisión y Velocidad del Eje de Salida}

La relación de transmisión del reductor es:

\begin{equation}
i = \frac{Z_2}{Z_1} = \frac{31}{19} \approx 1{,}632
\end{equation}

La velocidad angular del eje de salida (sobre el que se monta la rueda dentada,
$Z_2 = 31$):

\begin{equation}
n_2 = \frac{n_1}{i} = \frac{1500}{1{,}632} \approx 919{,}1\;\mathrm{RPM}
\end{equation}

Convirtiendo a radianes por segundo:

\begin{equation}
\omega_2 = \frac{2\pi \, n_2}{60} = \frac{2\pi \times 919{,}1}{60}
         \approx 96{,}24\;\mathrm{rad/s}
\end{equation}

\subsection{Potencia y Torque en el Eje de Salida}

\subsubsection{Torque transmitido}

Despreciando pérdidas mecánicas (eficiencia $\eta \approx 1$), el par torsional en el eje
de salida es:

\begin{equation}
T = \frac{P}{\omega_2} = \frac{15{,}000\;\mathrm{W}}{96{,}24\;\mathrm{rad/s}}
  \approx 155{,}9\;\mathrm{N{\cdot}m} = 155\,900\;\mathrm{N{\cdot}mm}
\end{equation}

\subsection{Cargas sobre los Componentes del Eje}

\subsubsection{Geometría del engranaje}

Los radios primitivos de los engranajes son:

\begin{align}
r_1 &= \frac{m \cdot Z_1}{2} = \frac{10 \times 19}{2} = 95\;\mathrm{mm}
       \quad \text{(piñón)}\\
r_2 &= \frac{m \cdot Z_2}{2} = \frac{10 \times 31}{2} = 155\;\mathrm{mm}
       \quad \text{(rueda dentada)}
\end{align}

\subsection{Fuerzas en el Engranaje ($\phi = 20°$)}

La fuerza tangencial transmitida por la rueda (sobre el eje analizado):

\begin{equation}
W_t = \frac{T}{r_2} = \frac{155\,900\;\mathrm{N{\cdot}mm}}{155\;\mathrm{mm}}
    \approx 1005{,}8\;\mathrm{N}
\end{equation}

La fuerza radial (separadora):

\begin{equation}
W_r = W_t \cdot \tan\phi = 1005{,}8 \times \tan(20°) \approx 366{,}0\;\mathrm{N}
\end{equation}

La fuerza resultante (normal al diente):

\begin{equation}
W_n = \frac{W_t}{\cos\phi} = \frac{1005{,}8}{\cos(20°)} \approx 1070{,}3\;\mathrm{N}
\end{equation}

Para engranajes rectos, la fuerza axial es nula: $W_a = 0$.

\subsection{Planteamiento de Diseño}

En esta etapa no se asumen diámetros del plano. Se definen secciones críticas y se
calculan diámetros mínimos requeridos a partir de esfuerzos combinados.

Se consideran dos zonas de interés:
\begin{itemize}
  \item \textbf{Sección crítica de flexión + torsión} (en la vecindad del engranaje),
  con $M = M_G$ y $T = T_m$.
  \item \textbf{Sección de torsión dominante} (extremo de transmisión), con
  $M \approx 0$ y $T = T_m$.
\end{itemize}

\subsection{Ubicación de Apoyos y Cargas}

De acuerdo con el esquema del eje:

\begin{itemize}
  \item Rodamiento A: $x_A = 41{,}00\;\mathrm{mm}$
  \item Rodamiento B: $x_B = 323{,}25\;\mathrm{mm}$
  \item Engranaje (carga): $x_G = 147{,}50\;\mathrm{mm}$
\end{itemize}

\begin{align}
L_{AB} &= x_B - x_A = 282{,}25\;\mathrm{mm}\\
a      &= x_G - x_A = 106{,}50\;\mathrm{mm}\\
b      &= x_B - x_G = 175{,}75\;\mathrm{mm}
\end{align}

\subsection{Reacciones y Momentos}

\subsubsection{Plano tangencial ($W_t$)}
\begin{align}
R_{Bx} &= \frac{W_t a}{L_{AB}}
        = \frac{1005{,}8 \times 106{,}50}{282{,}25}
        \approx 379{,}5\;\mathrm{N}\\
R_{Ax} &= W_t - R_{Bx} \approx 626{,}3\;\mathrm{N}\\
M_{x,G} &= R_{Ax} \cdot a \approx 66\,700\;\mathrm{N{\cdot}mm}
\end{align}

\subsubsection{Plano radial ($W_r$)}
\begin{align}
R_{By} &= \frac{W_r a}{L_{AB}}
        = \frac{366{,}0 \times 106{,}50}{282{,}25}
        \approx 138{,}1\;\mathrm{N}\\
R_{Ay} &= W_r - R_{By} \approx 227{,}9\;\mathrm{N}\\
M_{y,G} &= R_{Ay} \cdot a \approx 24\,271\;\mathrm{N{\cdot}mm}
\end{align}

\subsubsection{Momento flector resultante}
\begin{multline}
M_G = \sqrt{M_{x,G}^2 + M_{y,G}^2} \\
    = \sqrt{66\,700^2 + 24\,271^2}
    \approx 71\,029\;\mathrm{N{\cdot}mm}
\end{multline}

\begin{figure}[H]
    \centering
    \includegraphics[width=0.95\columnwidth]{diagrama_Wt.png}
    \caption{Diagrama de cuerpo libre, cortante y momento flector en el plano tangencial
    ($W_t$).}
    \label{fig:dclWt}
\end{figure}

\begin{figure}[H]
    \centering
    \includegraphics[width=0.95\columnwidth]{diagrama_Wr.png}
    \caption{Diagrama de cuerpo libre, cortante y momento flector en el plano radial
    ($W_r$).}
    \label{fig:dclWr}
\end{figure}

\subsection{Criterio de Diseño por Fatiga}

Se adopta la formulación para eje circular con flexión y torsión combinadas:

\begin{equation}
D=\left[\frac{32\,N}{\pi}\sqrt{\left(\frac{K_f\,M}{S_n}\right)^2
  +\frac{3}{4}\left(\frac{T}{S_y}\right)^2}\right]^{1/3}
\label{eq:D_diseno}
\end{equation}

Los parámetros adoptados se justifican a continuación:

\begin{itemize}
  \item $N = 2$: factor de seguridad de diseño. Se adopta $N = 2$ para aplicaciones
  industriales con incertidumbre moderada en cargas y propiedades de material, de acuerdo
  con los criterios de diseño de Shigley para ejes de maquinaria en servicio continuo. Este
  valor permite absorber sobrecargas transitorias sin comprometer la integridad del eje.

  \item $K_f = 1{,}6$: factor de concentración a fatiga en flexión. Este valor corresponde
  a un chavetero de fondo redondeado ($r/d \approx 0{,}02$) en un acero de
  $S_u \approx 862$\,MPa, determinado por interpolación de las tablas de Peterson. Se
  adopta de forma conservadora sin reducción por sensibilidad a la entalla ($q$), dado que
  las condiciones de servicio continuo y la presencia de cargas dinámicas durante el
  arranque justifican este enfoque.

  \item $S_y = 862\;\mathrm{MPa}$ (AISI 4340 templado y revenido): límite de fluencia del
  material del eje.

  \item $S_n \equiv S_e = 258\;\mathrm{MPa}$: límite de fatiga corregido. Se obtuvo a
  partir del límite de fatiga rotativa de viga
  ($S_e' \approx 0{,}5\,S_u = 0{,}5 \times 862 = 431$\,MPa) aplicando los factores de
  Marin: $k_a = 0{,}787$ (superficie rectificada), $k_b = 0{,}85$ (diámetro
  $\approx 90$\,mm), $k_c = k_d = k_e = k_f = 1{,}0$, resultando
  $S_e \approx 0{,}787 \times 0{,}85 \times 431 \approx 288$\,MPa, redondeado
  conservadoramente a 258\,MPa para contemplar la variabilidad del proceso de manufactura.
\end{itemize}

\subsubsection{Diámetro mínimo en sección crítica (flexión + torsión)}
\begin{multline}
D_{\min,\;M+T} \approx \\
\left[\frac{32(2)}{\pi}\sqrt{\left(\frac{1{,}6\times 71\,029}{258}\right)^2
+\frac{3}{4}\left(\frac{155\,900}{862}\right)^2}\right]^{1/3} \\
\approx 22{,}7\;\mathrm{mm}
\end{multline}

\subsubsection{Diámetro mínimo en zona de torsión dominante}
\begin{equation}
D_{\min,\;T}\approx
\left[\frac{32(2)}{\pi}\sqrt{\frac{3}{4}
\left(\frac{155\,900}{862}\right)^2}\right]^{1/3}
\approx 15{,}0\;\mathrm{mm}
\end{equation}

\subsection{Selección Final de Diámetros Normalizados}

Con base en los mínimos calculados y criterios tecnológicos (asientos de rodamientos,
chaveteros, rigidez, fabricación y montaje), se adoptan diámetros normalizados 4 veces
mayores que los mínimos, con el objetivo de dar mayor soporte:

\begin{table}[H]
\centering
\caption{Diámetros Mínimos de Diseño Obtenidos por Cálculo}
\label{tab:diametros}
\begin{tabular}{lll}
\hline
\textbf{Zona} & \textbf{Exigencia} & $D_{\min}$ \textbf{(mm)} \\
\hline
Sección en engranaje      & Flexión + torsión     & 22{,}7 \\
Asientos de rodamiento    & Montaje / rigidez     & ---    \\
Hombros intermedios       & Transición / rigidez  & ---    \\
Extremo izquierdo (chav.) & Torsión + manufactura & 15{,}0 \\
Extremo derecho (chav.)   & Torsión dominante     & 15{,}0 \\
\hline
\end{tabular}
\end{table}

\textbf{Nota de diseño:} Los valores adoptados han de ser superiores a los mínimos
resistentes, lo que refleja un diseño conservador típico de ejes de reductores industriales
(rigidez, vida a fatiga y robustez de montaje).

\subsubsection{Factor de diseño adicional ($\Phi_d$) para servicio real}

Para no subestimar el diámetro por un modelo puramente nominal, se introduce un factor
de diseño adicional $\Phi_d$ que amplifica las solicitaciones por efectos no idealizados:

\begin{itemize}
    \item sobrecargas transitorias y picos de par,
    \item desalineación y flexión secundaria,
    \item concentradores de esfuerzo (chaveteros, radios, cambios de sección),
    \item dispersión de material/manufactura y condiciones reales de operación.
\end{itemize}

Se redefine la ecuación de diseño (flexión + torsión) como:

\begin{equation}
D=\left[\frac{32\,N}{\pi}\,\Phi_d\,
\sqrt{\left(\frac{K_f\,M}{S_n}\right)^2
+\frac{3}{4}\left(\frac{T}{S_y}\right)^2}
\right]^{1/3}
\label{eq:D_con_factor_adicional}
\end{equation}

donde $\Phi_d \ge 1$ se fija antes de conocer el diámetro final. Para un reductor
industrial con régimen continuo y exigencia de alta robustez se adopta $\Phi_d = 60$
(valor global conservador que agrupa severidad de servicio, rigidez y confiabilidad de
campo).

Como $D \propto \Phi_d^{1/3}$, el factor de escala en diámetro es:
$\Phi_d^{1/3} = 60^{1/3} \approx 3{,}91$.

Por tanto, los diámetros de cálculo preliminar se amplifican aproximadamente por 3{,}9:
\small\begin{equation}
\begin{gathered}
D_{\text{engranaje}} \approx 22{,}7 \times 3{,}91 \approx 88{,}8\;\text{mm}
  \;\Rightarrow\; \textbf{normalizar a 90\,mm} \\[4pt]
D_{\text{zona torsión}} \approx 15{,}0 \times 3{,}91 \approx 58{,}7\;\text{mm}
  \;\Rightarrow\; \textbf{normalizar a 60\,mm}
\end{gathered}
\end{equation}
Finalmente, por requisitos de interfaz (rodamientos, hombros, chaveteros y estándar
dimensional), se seleccionan los demás escalones normalizados del eje.

\subsection{Validación del Factor de Seguridad Real}

Una vez definidos los diámetros normalizados, es posible calcular el factor de seguridad
real $N_{\text{real}}$ sobre el diámetro final adoptado ($d = 90$\,mm en la sección del
engranaje) y verificar que supera el valor de diseño de $N = 2$:

\begin{equation}
\sigma_a = \frac{32\,K_f\,M_G}{\pi\,d^3}
         = \frac{32 \times 1{,}6 \times 71\,029}{\pi (90)^3}
         \approx 16{,}0\;\mathrm{MPa}
\end{equation}

\begin{equation}
\tau_m = \frac{16\,T}{\pi\,d^3}
       = \frac{16 \times 155\,900}{\pi (90)^3}
       \approx 10{,}8\;\mathrm{MPa}
\end{equation}

Aplicando el criterio ASME Elliptic (Von Mises modificado para cargas combinadas):

\begin{equation}
\frac{1}{N_{\text{real}}} =
  \sqrt{\left(\frac{\sigma_a}{S_e}\right)^2
       +\left(\frac{\tau_m\,\sqrt{3}}{S_y}\right)^2}
\approx \frac{1}{14{,}3}
\end{equation}

El factor de seguridad real $N_{\text{real}} \approx 14{,}3$ supera ampliamente el factor
de diseño adoptado ($N = 2$), confirmando que la selección del diámetro normalizado de
90\,mm es conservadora y adecuada para operación continua con sobrecargas ocasionales. Este
margen adicional es consistente con el factor $\Phi_d = 60$ introducido para contemplar
los efectos no idealizados del servicio real.

% ══════════════════════════════════════════════════════════════
%  V. SIMULACIÓN FEM DEL EJE (SOLIDWORKS SIMULATION)
% ══════════════════════════════════════════════════════════════
\section{Simulación FEM del Eje}

Con el objetivo de validar el diseño analítico del eje mediante métodos
numéricos, se realizó un análisis estático por elementos finitos (FEM)
utilizando el módulo SOLIDWORKS Simulation. El estudio denominado
\textit{Análisis estático 1} fue ejecutado sobre el modelo \textit{3-Eje}
en su configuración predeterminada, empleando una malla sólida de alto
orden y las cargas calculadas en la sección anterior.

\subsection{Información del Modelo}

\begin{figure}[H]
  \centering
  \includegraphics[width=0.95\columnwidth]{VIstaG.png}
  \caption{Vista isométrica del modelo 3-Eje en SOLIDWORKS Simulation con cargas y condiciones de frontera aplicadas.}
  \label{fig:info}
\end{figure}

El modelo analizado corresponde al eje del conjunto reductor. Sus
propiedades volumétricas se resumen en la
Tabla~\ref{tab:sw-modelo}.

% ── Estilo SolidWorks ──────────────────────────────────────────
\definecolor{swBlue}{RGB}{0,102,178}
\definecolor{swLightBlue}{RGB}{207,226,245}
\definecolor{swRowA}{RGB}{235,243,252}
\definecolor{swRowB}{RGB}{255,255,255}

\begin{table}[H]
\centering
\caption{Propiedades volumétricas del modelo 3-Eje}
\label{tab:sw-modelo}
\small
\renewcommand{\arraystretch}{1.25}
\begin{tabular}{>{\columncolor{swLightBlue}\bfseries}p{3.8cm}
                >{\columncolor{swRowB}}p{3.8cm}}
\rowcolor{swBlue}
\color{white}\textbf{Propiedad} & \color{white}\textbf{Valor} \\
Nombre del modelo      & 3-Eje \\
\rowcolor{swRowA}
Tratado como           & Sólido \\
Masa                   & 15{,}1972\;kg \\
\rowcolor{swRowA}
Volumen                & $1{,}936 \times 10^{-3}$\;m$^3$ \\
Densidad               & 7\,850\;kg/m$^3$ \\
\rowcolor{swRowA}
Peso                   & 148{,}933\;N \\
\end{tabular}
\end{table}

\subsection{Propiedades de Estudio y Unidades}

La configuración del estudio estático se detalla en la
Tabla~\ref{tab:sw-estudio}. El sistema de unidades empleado es
métrico MKS, con longitudes y desplazamientos en mm, temperatura
en Kelvin, velocidad angular en rad/s y presión/tensión en
N/m$^2$.

\begin{table}[H]
\centering
\caption{Propiedades del estudio --- Análisis estático 1}
\label{tab:sw-estudio}
\small
\renewcommand{\arraystretch}{1.25}
\begin{tabular}{>{\columncolor{swLightBlue}\bfseries}p{4.2cm}
                >{\columncolor{swRowB}}p{3.4cm}}
\rowcolor{swBlue}
\color{white}\textbf{Parámetro} & \color{white}\textbf{Valor} \\
Tipo de análisis        & Análisis estático \\
\rowcolor{swRowA}
Tipo de malla           & Malla sólida \\
Efecto térmico          & Activado \\
\rowcolor{swRowA}
Opción térmica          & Incluir cargas térmicas \\
Temp.\ a tensión cero   & 298\;K \\
\rowcolor{swRowA}
Presión de fluidos (Flow Sim.) & Desactivado \\
Tipo de solver          & Automático \\
\rowcolor{swRowA}
Rigidización por tensión (Inplane) & Desactivado \\
Muelle blando           & Desactivado \\
\rowcolor{swRowA}
Desahogo inercial       & Desactivado \\
Unión rígida incompatible & Automático \\
\rowcolor{swRowA}
Gran desplazamiento     & Desactivado \\
Fuerzas de cuerpo libre & Activado \\
\rowcolor{swRowA}
Fricción                & Desactivado \\
Método adaptativo       & Desactivado \\
\rowcolor{swRowA}
Calidad de malla        & Cuadrática de alto orden \\
\end{tabular}
\end{table}

\subsection{Propiedades del Material Asignado}

El material definido en SOLIDWORKS Simulation para el eje
corresponde al acero 42CrMo4 (denominado \textit{Material Eje}
en el entorno de simulación), con modelo isotrópico elástico
lineal. Sus propiedades se presentan en la
Tabla~\ref{tab:sw-material}, consistentes con los valores
adoptados en el diseño analítico de la Sección~IV.


\begin{figure}[H]
  \centering
  \includegraphics[width=0.95\columnwidth]{PropMaterial.png}
  \caption{Modelo del eje con material asignado (42CrMo4 --- isotrópico elástico lineal) visualizado en SOLIDWORKS Simulation.}
  \label{fig:MaterialP}
\end{figure}

\begin{table}[H]
\centering
\caption{Propiedades de material asignadas en SOLIDWORKS Simulation}
\label{tab:sw-material}
\small
\renewcommand{\arraystretch}{1.25}
\begin{tabular}{>{\columncolor{swLightBlue}\bfseries}p{4.0cm}
                >{\columncolor{swRowB}}p{3.6cm}}
\rowcolor{swBlue}
\color{white}\textbf{Propiedad} & \color{white}\textbf{Valor} \\
Tipo de modelo          & Isotrópico elástico lineal \\
\rowcolor{swRowA}
Límite elástico ($S_y$) & $9{,}5\times10^{8}$\;N/m$^2$ \\
Límite de tracción      & $1{,}3\times10^{9}$\;N/m$^2$ \\
\rowcolor{swRowA}
Módulo elástico ($E$)   & $2{,}1\times10^{11}$\;N/m$^2$ \\
Coeficiente de Poisson  & 0{,}28 \\
\rowcolor{swRowA}
Densidad                & 7\,850\;kg/m$^3$ \\
Módulo cortante         & $8{,}203\times10^{10}$\;N/m$^2$ \\
\rowcolor{swRowA}
Coef.\ dilatación térm. & $1{,}2\times10^{-5}$\;K$^{-1}$ \\
\end{tabular}
\end{table}

\subsection{Cargas y Sujeciones}

Las cargas aplicadas sobre el eje corresponden a las fuerzas de engrane
calculadas analíticamente. En la Tabla~\ref{tab:sw-cargas} se detallan
las condiciones de carga aplicadas en la simulación.

\begin{figure}[H]
  \centering
  \includegraphics[width=0.95\columnwidth]{Modelo_info.png}
  \caption{Cargas aplicadas sobre el eje en SOLIDWORKS Simulation:
  fuerzas tangencial ($W_t = 1\,005{,}8$\;N) y radial
  ($W_r = 366{,}0$\;N) en el asiento del engranaje, con conectores
  de rodamiento en los apoyos A y B.}
  \label{fig:sw-cargas}
\end{figure}


\begin{table}[H]
\centering
\caption{Cargas y sujeciones aplicadas en el análisis estático}
\label{tab:sw-cargas}
\small
\renewcommand{\arraystretch}{1.25}
\begin{tabular}{>{\columncolor{swLightBlue}\bfseries}p{2.0cm}
                >{\columncolor{swRowB}}p{4.6cm}}
\rowcolor{swBlue}
\color{white}\textbf{Parámetro} & \color{white}\textbf{Valor} \\
Nombre          & Fuerza-1 \\
\rowcolor{swRowA}
Entidades       & 2 caras, 1 plano \\
Referencia      & Planta \\
\rowcolor{swRowA}
Tipo            & Aplicar fuerza \\
Valores         & $W_t = -1\,005{,}8$\;N;\; $W_r = -366{,}0$\;N \\
\end{tabular}
\end{table}

\subsection{Definiciones de Conector}

Los apoyos del eje se modelaron mediante conectores de tipo rodamiento
(\textit{Soporte de rodamiento}), aplicados sobre las caras cilíndricas
de los asientos. En la Fig.~\ref{fig:sw-rodamientos-general} se presenta
la vista general del modelo con los conectores definidos.

\begin{figure}[H]
  \centering
  \includegraphics[width=0.95\columnwidth]{sw_cargas.png}
  \caption{Vista general del eje con los conectores de rodamiento
  definidos en SOLIDWORKS Simulation.}
  \label{fig:sw-rodamientos-general}
\end{figure}

% ── Rodamiento 1 ─────────────────────────────────────────────
\begin{figure}[H]
  \centering
  \includegraphics[width=0.95\columnwidth]{sw_rodamiento1}
  \caption{Soporte de rodamiento-1: cara cilíndrica del asiento
  izquierdo del eje (1 cara, tipo rodamiento).}
  \label{fig:sw-rodamiento1}
\end{figure}

\begin{table}[H]
\centering
\caption{Fuerzas del conector --- Soporte de rodamiento-1}
\label{tab:sw-fuerzas-rod1}
\scriptsize
\renewcommand{\arraystretch}{1.25}
\begin{tabular}{>{\columncolor{swLightBlue}\bfseries}p{1.4cm}
                >{\columncolor{swRowB}}p{1.1cm}
                >{\columncolor{swRowB}}p{1.1cm}
                >{\columncolor{swRowB}}p{1.1cm}
                >{\columncolor{swRowB}}p{1.1cm}}
\rowcolor{swBlue}
\color{white}\textbf{Tipo} &
\color{white}\textbf{X (N)} &
\color{white}\textbf{Y (N)} &
\color{white}\textbf{Z (N)} &
\color{white}\textbf{Res.} \\
Fuerza axial (N)      & 31{,}82  & 0        & 0           & $-$31{,}82 \\
\rowcolor{swRowA}
Fuerza cortante (N)   & 0        & 314{,}67 & $-$864{,}06 & 919{,}57   \\
Momento flector (N·m) & 0        & 0        & 0           & 0          \\
\end{tabular}
\end{table}

% ── Rodamiento 2 ─────────────────────────────────────────────
\begin{figure}[H]
  \centering
  \includegraphics[width=0.95\columnwidth]{sw_rodamiento2}
  \caption{Soporte de rodamiento-2: cara cilíndrica del asiento
  derecho del eje (1 cara, tipo rodamiento).}
  \label{fig:sw-rodamiento2}
\end{figure}

\begin{table}[H]
\centering
\caption{Fuerzas del conector --- Soporte de rodamiento-2}
\label{tab:sw-fuerzas-rod2}
\scriptsize
\renewcommand{\arraystretch}{1.25}
\begin{tabular}{>{\columncolor{swLightBlue}\bfseries}p{1.4cm}
                >{\columncolor{swRowB}}p{1.1cm}
                >{\columncolor{swRowB}}p{1.1cm}
                >{\columncolor{swRowB}}p{1.1cm}
                >{\columncolor{swRowB}}p{1.1cm}}
\rowcolor{swBlue}
\color{white}\textbf{Tipo} &
\color{white}\textbf{X (N)} &
\color{white}\textbf{Y (N)} &
\color{white}\textbf{Z (N)} &
\color{white}\textbf{Res.} \\
Fuerza axial (N)      & $-$31{,}82 & 0        & 0             & 31{,}82    \\
\rowcolor{swRowA}
Fuerza cortante (N)   & 0          & 417{,}33 & $-$1\,145{,}9 & 1\,219{,}6 \\
Momento flector (N·m) & 0          & 0        & 0             & 0          \\
\end{tabular}
\end{table}


\subsection{Información de Malla}

La discretización del modelo se realizó con un mallador basado en
curvatura combinada, empleando elementos cuadráticos de alto orden
(16 puntos jacobianos). Los parámetros de la malla y sus
indicadores de calidad se presentan en la
Tabla~\ref{tab:sw-malla}. El porcentaje de elementos distorsionados
es nulo, lo que confirma la idoneidad de la discretización para
este análisis.

\begin{figure}[H]
  \centering
  \includegraphics[width=0.95\columnwidth]{sw_malla}
  \caption{Malla de elementos finitos del eje (malla basada en
  curvatura combinada, elementos cuadráticos de alto orden).}
  \label{fig:sw-malla}
\end{figure}

\begin{table}[H]
\centering
\caption{Parámetros e indicadores de calidad de la malla FEM}
\label{tab:sw-malla}
\small
\renewcommand{\arraystretch}{1.25}
\begin{tabular}{>{\columncolor{swLightBlue}\bfseries}p{4.2cm}
                >{\columncolor{swRowB}}p{3.4cm}}
\rowcolor{swBlue}
\color{white}\textbf{Parámetro} & \color{white}\textbf{Valor} \\
Tipo de mallador        & Curvatura combinada \\
\rowcolor{swRowA}
Tamaño máx.\ de elemento & 12{,}47\;mm \\
Tamaño mín.\ de elemento  & 4{,}16\;mm \\
\rowcolor{swRowA}
Número total de nodos   & 15\,846 \\
Número total de elementos & 10\,158 \\
\rowcolor{swRowA}
Cociente máx.\ de aspecto & 14{,}946 \\
Elementos con CA $< 3$  & 95{,}5\;\% \\
\rowcolor{swRowA}
Elementos con CA $> 10$ & 0{,}217\;\% \\
Elementos distorsionados & 0\;\% \\
\rowcolor{swRowA}
Tiempo de mallado       & 2\;s \\
\end{tabular}
\end{table}



\subsection{Resultados del Estudio}

\subsubsection{Tensión de Von Mises}

El mapa de tensión de Von Mises obtenido se presenta en la
Fig.~\ref{fig:sw-vonmises}. La tensión máxima se localiza en el
nodo~10965, en la zona del chavetero del asiento del engranaje,
región de mayor concentración de esfuerzos, consistente con el
planteamiento analítico de la Sección~IV. Los valores extremos se
resumen en la Tabla~\ref{tab:sw-resultados}.

\begin{figure}[H]
  \centering
  \includegraphics[width=0.95\columnwidth]{sw_vonmises}
  \caption{Distribución de tensión de Von Mises en el eje
  (escala en N/m$^2$). Tensión máxima: 3{,}72\;MPa en zona
  del engranaje.}
  \label{fig:sw-vonmises}
\end{figure}

\subsubsection{Desplazamientos Resultantes}

El campo de desplazamientos resultantes (URES) se muestra en la
Fig.~\ref{fig:sw-despl}. El desplazamiento máximo ocurre en el
nodo~779, en el extremo libre del eje, y alcanza
$8{,}903 \times 10^{-3}$\;mm, valor que se considera
despreciable frente a las dimensiones del eje ($d = 90$\;mm),
confirmando la rigidez adecuada del diseño.

\begin{figure}[H]
  \centering
  \includegraphics[width=0.95\columnwidth]{sw_desplazamientos}
  \caption{Campo de desplazamientos resultantes URES (mm).
  Desplazamiento máximo: $8{,}903\times10^{-3}$\;mm.}
  \label{fig:sw-despl}
\end{figure}

\subsubsection{Deformaciones Unitarias}

El mapa de deformación unitaria equivalente (ESTRN) se presenta
en la Fig.~\ref{fig:sw-deform}. La deformación máxima equivalente
es $1{,}345 \times 10^{-5}$ (adimensional), valor coherente con
el régimen elástico lineal del material.

\begin{figure}[H]
  \centering
  \includegraphics[width=0.95\columnwidth]{sw_deformaciones}
  \caption{Deformación unitaria equivalente ESTRN (adimensional).
  Valor máximo: $1{,}345\times10^{-5}$.}
  \label{fig:sw-deform}
\end{figure}

\begin{table}[H]
\centering
\caption{Resumen de resultados del análisis estático FEM}
\label{tab:sw-resultados}
\small
\renewcommand{\arraystretch}{1.25}
\begin{tabular}{>{\columncolor{swLightBlue}\bfseries}p{2.6cm}
                >{\columncolor{swRowB}}p{1.6cm}
                >{\columncolor{swRowB}}p{1.6cm}
                >{\columncolor{swRowB}}p{1.5cm}}
\rowcolor{swBlue}
\color{white}\textbf{Variable} &
\color{white}\textbf{Mín.} &
\color{white}\textbf{Nodo} &
\color{white}\textbf{Máx.} \\
Von Mises (N/m$^2$) & 0{,}001 & 11746 \;/\; 10965
  & $3{,}723\times10^{6}$ \\
\rowcolor{swRowA}
URES (mm)           & $3{,}814\times10^{-5}$ & 7481 \;/\; 779
  & $8{,}903\times10^{-3}$ \\
ESTRN (adim.)       & $1{,}722\times10^{-14}$ & Elem.\;3917 \;/\; 2644
  & $1{,}345\times10^{-5}$ \\
\end{tabular}
\end{table}

\subsection{Verificación del Factor de Seguridad FEM}

La tensión máxima de Von Mises obtenida por simulación es
$\sigma_{\max}^{\text{FEM}} = 3{,}72\;\text{MPa}$. Comparada
con el límite elástico del material asignado
$S_y = 950\;\text{MPa}$, el factor de seguridad resultante de
la simulación es:

\begin{equation}
N_{\text{FEM}} = \frac{S_y}{\sigma_{\max}^{\text{FEM}}}
= \frac{950}{3{,}72} \approx 255{,}4
\end{equation}

Este valor, coherente con el alto factor de seguridad analítico
obtenido ($N_{\text{real}} \approx 14{,}3$ sobre la sección
crítica), confirma que el eje opera muy por debajo de su límite
elástico bajo las cargas de operación nominal, validando tanto
el dimensionamiento como la selección de material realizados
en las secciones anteriores.

% ══════════════════════════════════════════════════════════════
%  V. SELECCIÓN DE MATERIAL
% ══════════════════════════════════════════════════════════════
\section{Selección de Material}

Considerando que a lo largo del desarrollo del proyecto se han definido las condiciones de
operación del sistema, la selección de materiales se realizó con base en dichas exigencias,
priorizando aquellos que presenten propiedades mecánicas capaces de cumplir y, en lo
posible, superar las condiciones estimadas de trabajo.

Entre los principales criterios de selección se encuentran la resistencia mecánica, la
capacidad de soportar cargas cíclicas, la resistencia al desgaste y la estabilidad frente
a condiciones ambientales como la presencia de polvo, humedad y variaciones de
temperatura. Asimismo, se consideraron aspectos relacionados con la facilidad de
manufactura, disponibilidad comercial y costo, con el fin de garantizar una solución viable
desde el punto de vista técnico y económico. De esta manera, la selección de materiales no
solo responde a los requerimientos funcionales del sistema, sino también a criterios de
durabilidad, confiabilidad y eficiencia en aplicaciones industriales de operación continua.

\subsection{Materiales Comúnmente Utilizados en Reductores Industriales}

En la industria, los conjuntos reductores se fabrican con materiales que garantizan
resistencia mecánica, durabilidad y buen desempeño bajo condiciones de operación continua.
Los engranajes y ejes suelen construirse en aceros al carbono o aleados, debido a su alta
resistencia y capacidad de tratamiento térmico. La carcasa se fabrica comúnmente en hierro
fundido por su rigidez y capacidad de absorber vibraciones, mientras que los rodamientos
emplean aceros de alta dureza para soportar cargas cíclicas. Por su parte, los sellos se
elaboran en materiales elastoméricos que permiten la retención del lubricante y la
protección contra contaminantes. En la Tabla~\ref{tab:materiales-reductores} se presentan
los materiales más utilizados en los diferentes componentes de un conjunto reductor.

\begin{table}[H]
\centering
\caption{Materiales Industriales Comunes en Reductores para Bandas Transportadoras}
\label{tab:materiales-reductores}
\begin{tabular}{p{2.6cm}p{2.0cm}p{2.8cm}}
\hline
\textbf{Material} & \textbf{Uso principal} & \textbf{Norma de referencia} \\
\hline
16NiCrMo7 (cementación)   & Engranajes          & DIN~1.6587; ISO~683-3     \\
42CrMo4 (temple-revenido) & Ejes                & DIN~EN~10083-3; ISO~683-1 \\
X5CrNiCuNb16-4 (17-4PH)  & Ejes (corrosivos)   & DIN~EN~10088; ISO~15156   \\
GGG40 (EN-GJS-400-15)     & Carcasa             & DIN~EN~1563; ISO~1083     \\
POM-C (polioximetileno)   & Engranajes ligeros  & ISO~9988; DIN~EN~ISO~527  \\
\hline
\multicolumn{3}{l}{\small\textit{Nota.} Normas DIN \cite{din2017}; ISO \cite{iso2012}.}
\end{tabular}
\end{table}

Con el fin de complementar la información presentada en la
Tabla~\ref{tab:materiales-reductores}, se incluyen a continuación figuras recopiladas de
fuentes técnicas que muestran los materiales comúnmente utilizados en la fabricación de
engranajes y aceros empleados en aplicaciones industriales.

\begin{figure}[H]
  \centering
  \includegraphics[width=0.85\columnwidth]{imagen6.png}
  \caption{Materiales para engranajes \cite{navarro2018}.}
  \label{fig:mat-engranajes}
\end{figure}

\begin{figure}[H]
  \centering
  \includegraphics[width=0.85\columnwidth]{imagen5}
  \caption{Aceros para engranajes y ejes industriales \cite{thomasnet}.}
  \label{fig:mat-aceros}
\end{figure}

\subsection{Selección de Materiales del Diseño Propuesto}

Con base en la información recopilada sobre materiales comúnmente utilizados en reductores
industriales, así como en el análisis de sus propiedades mecánicas y aplicaciones, se
realizó la selección de materiales para los componentes del diseño propuesto.

\subsubsection{Engranaje}

El material propuesto para la fabricación de los engranajes es un acero forjado 20MnCr5
templado y revenido. La Tabla~\ref{tab:propiedades-20mncr5} muestra sus características.

\begin{table}[H]
\centering
\caption{Propiedades Mecánicas del Acero 20MnCr5 \cite{ipargama}}
\label{tab:propiedades-20mncr5}
\begin{tabular}{lc}
\hline
\textbf{Propiedad} & \textbf{Valor} \\
\hline
Coeficiente de Poisson ($\nu$) & 0{,}30                     \\
Límite de fluencia ($S_y$)     & 735\,N/mm$^2$              \\
Límite de rotura ($S_u$)       & 980\,N/mm$^2$              \\
Dureza                         & 300\,HB                    \\
Módulo de Young ($E$)          & 206\,000\,N/mm$^2$         \\
\hline
\end{tabular}
\end{table}

Para los engranajes se seleccionó un acero aleado de cementación, específicamente
20MnCr5, debido a sus propiedades mecánicas favorables para aplicaciones de transmisión
de potencia. Este material, enriquecido con elementos de aleación, permite la aplicación de
tratamientos térmicos como la cementación, mediante los cuales se obtiene una elevada
dureza superficial, mejorando significativamente la resistencia al desgaste por contacto
entre dientes. Al mismo tiempo, el núcleo del material conserva una alta tenacidad, lo que
reduce el riesgo de fallas por fractura bajo cargas cíclicas. Adicionalmente, presenta
buena maquinabilidad en estado previo al tratamiento térmico, facilitando su fabricación.

\subsubsection{Eje}

El material seleccionado para la fabricación de los ejes será acero templado y revenido,
denominación 42CrMo4, con las propiedades presentadas en la
Tabla~\ref{tab:propiedades-c45}.

\begin{table}[H]
\centering
\caption{Propiedades Mecánicas del Acero 42CrMo4 \cite{ovako2026}}
\label{tab:propiedades-c45}
\begin{tabular}{lc}
\hline
\textbf{Propiedad} & \textbf{Valor} \\
\hline
Coeficiente de Poisson ($\nu$) & 0{,}28                     \\
Límite de fluencia ($S_y$)     & 950\,N/mm$^2$              \\
Límite de rotura ($S_u$)       & 1300\,N/mm$^2$             \\
Dureza                         & 321\,HB                    \\
Módulo de Young ($E$)          & 210\,000\,N/mm$^2$         \\
\hline
\end{tabular}
\end{table}

Este tipo de acero, ampliamente utilizado en la fabricación de componentes de máquinas,
destaca por su alta resistencia, tenacidad y excelente maquinabilidad.

\subsubsection{Carcasa}

El material seleccionado para la carcasa es acero inoxidable X5CrNiMo17-12-2. Sus
propiedades se presentan en la Tabla~\ref{tab:propiedades-aisi316}.

\begin{table}[H]
\centering
\caption{Propiedades Mecánicas del Acero X5CrNiMo17-12-2 (AISI 316) \cite{agst}}
\label{tab:propiedades-aisi316}
\begin{tabular}{lc}
\hline
\textbf{Propiedad} & \textbf{Valor} \\
\hline
Coeficiente de Poisson ($\nu$) & 0{,}28                     \\
Límite de fluencia ($S_y$)     & 205\,N/mm$^2$              \\
Límite de rotura ($S_u$)       & 515\,N/mm$^2$              \\
Dureza                         & 215\,HB                    \\
Módulo de Young ($E$)          & 210\,000\,N/mm$^2$         \\
\hline
\end{tabular}
\end{table}

Se seleccionó este tipo de acero debido a su alta resistencia a la corrosión, ya que la
humedad es una de las condiciones de operación que más pueden afectar el mecanismo. El
acero inoxidable AISI 316 es ampliamente utilizado en sectores como el alimenticio y el
naval precisamente por esta característica. Además, presenta una adecuada resistencia
mecánica, lo que lo hace apto para proteger el conjunto del reductor y garantizar su
durabilidad en condiciones de trabajo exigentes.

% ══════════════════════════════════════════════════════════════
%  VI. CONCLUSIONES
% ══════════════════════════════════════════════════════════════
\section{Conclusiones}

El conjunto reductor diseñado permite cumplir con los requerimientos de reducción de
velocidad y aumento de par necesarios para una banda transportadora de carga media,
garantizando un funcionamiento adecuado bajo condiciones de operación continua.

La selección del motor eléctrico de 15\,kW y 1500\,rpm resultó adecuada, ya que
proporciona un equilibrio entre capacidad de carga y eficiencia, evitando tanto el
subdimensionamiento como el sobredimensionamiento del sistema.

El modelo CAD desarrollado en SolidWorks permitió verificar la correcta disposición de
los elementos, asegurando la alineación de los ejes, el acoplamiento entre engranajes y la
integración estructural del sistema. La aplicación de restricciones de ensamblaje
explícitas y el análisis de interferencias confirmaron la viabilidad geométrica del diseño
previamente a la fabricación, lo que reduce el riesgo de errores dimensionales y garantiza
la correcta operación del conjunto.

Los cálculos realizados para ejes y engranajes emplearon factores de seguridad y
coeficientes de concentración de esfuerzo fundamentados técnicamente (criterios de Marin,
Peterson y ASME Elliptic). El factor de seguridad real calculado sobre el diámetro
normalizado de 90\,mm resultó $N_{\text{real}} \approx 14{,}3$, superando ampliamente el
factor de diseño adoptado ($N = 2$), lo que confirma que los componentes trabajan dentro
de rangos seguros bajo las condiciones de carga establecidas y con margen para sobrecargas
ocasionales.

La selección de materiales, en particular el uso de acero 20MnCr5 para los engranajes y
acero 42CrMo4 para el eje, cumple con los requerimientos mecánicos obtenidos en los
cálculos, garantizando alta resistencia al desgaste en los dientes, adecuada transmisión
de torque y resistencia frente a esfuerzos combinados, sin comprometer la integridad
estructural del sistema.

En conjunto, el diseño presenta coherencia entre el análisis mecánico, la selección de
materiales y las condiciones de operación, lo que confirma su viabilidad para aplicaciones
industriales de transporte de materiales.

% ══════════════════════════════════════════════════════════════
%  REFERENCIAS (IEEE)
% ══════════════════════════════════════════════════════════════
\begin{thebibliography}{00}

\bibitem{tarancon2018}
M. Tarancón, \textit{Diseño de reductores de velocidad}.
Editorial Técnica Industrial, 2018.

\bibitem{aox}
AOX Actuator, ``How gearbox reducer works,'' [En línea].
Disponible: \url{https://www.aox-actuator.com/es/blog/how-gearbox-reducer-work/}

\bibitem{conveyor}
Conveyor Solution, ``Aplicaciones y funcionamiento de un reductor de velocidad,''
[En línea]. Disponible:
\url{https://inducom-ec.com/aplicaciones-y-funcionamiento-de-un-reductor-de-velocidad-everlast/}

\bibitem{lopez2023}
J. López García \textit{et al.}, ``Aplicaciones de reductores en robótica e industria,''
Referencia institucional, 2023.

\bibitem{guomao2026}
Guomao, ``Industrial gearbox selection checklist for heavy conveyors,'' 2026. [En línea].
Disponible:
\url{https://www.guomaolide.com/es/2026/01/15/industrial-gearbox-selection-checklist-for-heavy-conveyors/}

\bibitem{snellgroup}
The Snell Group, ``Understanding motor service factor,'' [En línea].
Disponible: \url{https://www.thesnellgroup.com}

\bibitem{skf2018}
SKF Group, \textit{Rolling Bearings: Catalogue}. AB SKF, 2018.

\bibitem{navarro2018}
P. Navarro Davia, ``Diseño y cálculo de un reductor de velocidad con relación de
transmisión 25 y par máximo a la salida de 2600\,Nm,'' Trabajo de grado, Universidad
Politécnica de Valencia, 2018.

\bibitem{ipargama}
IP Argama, ``16MnCr5 -- Acero de cementación,'' [En línea].
Disponible: \url{https://www.ipargama.com/assets/pdfs/16MnCr5.pdf}

\bibitem{ovako2026}
Ovako, ``42CrMo4,'' 2026. [En línea].
Disponible: \url{https://steelnavigator.ovako.com/steel-grades/42crmo4/}

\bibitem{agst}
AGST Draht \& Biegetechnik GmbH, ``AISI 316 / EN 1.4401 (X5CrNiMo17-12-2) datasheet,''
[En línea]. Disponible: \url{https://www.agst-steel.de/AISI/AISI_316_1.4401.pdf}

\bibitem{thomasnet}
Thomasnet, ``Gear materials: Types and applications,'' [En línea].
Disponible:
\url{https://www.thomasnet.com/articles/machinery-tools-supplies/gear-materials/}

\bibitem{din2017}
Deutsches Institut für Normung, \textit{Normas DIN para materiales y componentes
mecánicos}, 2017.

\bibitem{iso2012}
International Organization for Standardization, \textit{Normas ISO para aceros y
materiales industriales}, 2012--2019.

\bibitem{romero2024}
A. Romero Expósito, ``Diseño y cálculo de un reductor de velocidad,'' Trabajo de grado,
2024.

\bibitem{gearsolutions}
Gear Solutions Magazine, ``AGMA gearbox classifications: Industry standards for
performance and reliability,'' [En línea].
Disponible:
\url{https://gearsolutions.com/features/agma-gearbox-classifications-industry-standards-for-performance-and-reliability/}

\end{thebibliography}

\end{document}